\section*{Expected Outcomes}
This research primarily contributes theoretical models, analytical insights, and empirical validation of theory-driven claims. Implementation components serve as validation mechanisms for the proposed formulations rather than as primary contributions.

Upon completion, this PhD research is expected to deliver the following outcomes aligned with RQ1--RQ2:
\begin{itemize}
  \item \textbf{Joint optimization formulation (RQ1):} A formal combinatorial optimization model under a bounded-$K$ partition family that jointly determines partitioning and UAV assignment while minimizing end-to-end latency and total energy under coupled constraints.

  \item \textbf{Analytical characterization of $K$ (RQ1):} Structural analysis of how the number of partitions influences feasibility, complexity, and the latency--energy trade-off surface, including identification of regimes where communication overhead dominates and further partitioning yields diminishing returns.

  \item \textbf{Communication-aware cost model (RQ1):} An analytical model linking intermediate tensor size and time-varying FANET link conditions to transmission delay and RF energy within the optimization objective.

  \item \textbf{Online adaptation strategy (RQ2):} A characterization of how topology change, mobility, and failures alter the feasible region and objective landscape, together with a re-optimization mechanism that preserves near-optimal performance and inference continuity.

  \item \textbf{Empirical validation:} Co-simulation results (NS-3, ROS, Gazebo) validating the proposed formulations and adaptation logic across stress scenarios including node dropouts, link disruption, and varying mobility conditions.
\end{itemize}

Quantitative evaluation will assess the theoretical trade-off structure and adaptation behavior across varying system conditions.

Although instantiated for UAV swarms, the three-layer formulation (computation, network, control) is not aerial-specific: the same coupled structure arises in any mobile edge collective where compute placement is tied to a volatile communication substrate, such as ground robot fleets or maritime surface vehicles. UAV swarms constitute its most demanding instantiation, so results are expected to transfer downward to less constrained settings.
